\documentclass[11pt]{article}

\usepackage{nmrdoc}
\usepackage{siunitx}
\usepackage{bm}

\sisetup{per-mode=symbol}
\DeclareSIUnit{\angstrom}{\text{\AA}}
\DeclareSIUnit{\ppm}{ppm}

\begin{document}
\NmrTitle{The Biot--Savart Ring-Current Calculator}

\section{What the calculator computes}

The Biot--Savart calculator describes the magnetic field generated by an
aromatic ring current, using the Johnson--Bovey two-loop model of the
\(\pi\)-electron circulation. Aromatic electrons respond to an applied magnetic
field by circulating around the ring. That circulation produces a secondary
magnetic field at a nearby nucleus, and the secondary field changes the nuclear
shielding. The observable connected to this calculation is the NMR chemical
shift through its dependence on the shielding tensor.

The implementation does not solve for the electronic current from quantum
mechanics. It builds a geometry-derived tensor kernel from ring vertices, ring
normals, and fixed ring-type parameters. The file names use
\texttt{shielding}, but the code path writes the unit-current geometric kernel.
Calibration against a DFT shielding reference, or multiplication by an
explicit ring-current coefficient outside this routine, is what turns the
kernel into a shielding contribution in \(\si{\ppm}\).

\section{Where shielding enters}

For a nucleus in an applied field \(\bm B_0\), the local field is written
\begin{equation}
  \bm B_{\mathrm{nuc}}=(\bm I-\bm\sigma)\bm B_0 .
  \label{eq:local-field}
\end{equation}
\(\bm I\) is the \(3\times3\) identity matrix. The shielding tensor
\(\bm\sigma\) is dimensionless: its \(ab\) entry says how a field applied along
Cartesian direction \(b\) changes the field component felt by the nucleus along
direction \(a\). In component form,
\begin{equation}
  B^{\mathrm{sec}}_a=-\sum_b \sigma_{ab} B_{0,b}.
  \label{eq:secondary-field}
\end{equation}
The minus sign is the NMR shielding convention: positive shielding reduces the
field felt by the nucleus.

The quantum origin of \(\bm\sigma\) is electronic linear response. The applied
field perturbs the molecular electrons, the perturbed electrons create their
own magnetic field at the nucleus, and Eq.~\eqref{eq:secondary-field} records
the part proportional to \(\bm B_0\). A scalar is not enough because a molecule
has orientation. Turning the molecule with respect to \(\bm B_0\) can change
which component of the electronic response is excited and which component of
the secondary field arrives at the nucleus.

For an aromatic ring, this calculator keeps the part of that response that can
be represented as a classical current loop. If \(\hat{\bm n}\) is the unit
normal to the ring plane, then the model assumes that the ring current is
driven by the normal component \(\hat{\bm n}\cdot\bm B_0\) of the applied
field. The field produced by that current is then evaluated at the target
atom. This is an exact two-stage linear map within the model: first a
one-dimensional meter reads the projection of \(\bm B_0\) onto
\(\hat{\bm n}\); then that single reading drives a fixed three-component field
pattern at the atom. The \(3\times3\) tensor is the matrix of that projection
followed by that field pattern.

\section{The tensor split used by the code}

Any real \(3\times3\) tensor \(\bm X\) is decomposed by
\texttt{SphericalTensor::Decompose} into \(T_0+T_1+T_2\). The scalar part is
\begin{equation}
  T_0=\frac{1}{3}\operatorname{tr}\bm X .
  \label{eq:t0}
\end{equation}
It is the isotropic component. In rapidly tumbling solution, this trace
average is the part that enters the usual isotropic chemical shift after
reference subtraction.

The antisymmetric part is stored as the three-component vector
\begin{equation}
  \bm T_1=\frac{1}{2}
  \begin{pmatrix}
    X_{yz}-X_{zy}\\
    X_{zx}-X_{xz}\\
    X_{xy}-X_{yx}
  \end{pmatrix}.
  \label{eq:t1}
\end{equation}
This vector is the Levi-Civita dual of the antisymmetric matrix. It is usually
not observed directly in standard isotropic NMR, but it is present whenever
the response matrix is not symmetric.

The remaining part is the symmetric traceless matrix
\begin{equation}
  \bm S=\frac{\bm X+\bm X^{\mathsf T}}{2}-T_0\bm I,
  \qquad \operatorname{tr}\bm S=0 .
  \label{eq:sym-trace}
\end{equation}
It has five independent components. The code stores those five entries in an
isometric real spherical basis,
\begin{equation}
  \bm T_2 =
  \begin{pmatrix}
    \sqrt{2}\,S_{xy}\\
    \sqrt{2}\,S_{yz}\\
    \sqrt{3/2}\,S_{zz}\\
    \sqrt{2}\,S_{xz}\\
    (S_{xx}-S_{yy})/\sqrt{2}
  \end{pmatrix}.
  \label{eq:t2}
\end{equation}
``Isometric'' means that the five-vector length equals the Frobenius length of
the underlying symmetric traceless matrix:
\(\sum_m T_{2m}^2=\sum_{ab}S_{ab}^2\). The calculator writes all three parts in
the packed order
\[
  [T_0,T_{1x},T_{1y},T_{1z},
    T_{2,-2},T_{2,-1},T_{2,0},T_{2,+1},T_{2,+2}] .
\]
The Biot--Savart kernel is rank one, but it is generally not symmetric, so
\(T_0\), \(T_1\), and \(T_2\) can all be nonzero.

\section{The method implemented}

Each aromatic ring supplies a vertex list \(\bm v_i\) in \(\si{\angstrom}\), a
center, an average radius, and a unit normal \(\hat{\bm n}\). The center is the
mean of the vertices. The normal is the right singular vector of the centered
vertex matrix with the smallest singular value, with its sign chosen to match
the cyclic vertex order. This is the least-squares plane normal: it points
perpendicular to the least-squares plane through the ring atoms.

The finite-wire field used for each polygon edge is
\begin{equation}
  \bm F(\bm a,\bm b,I,\bm r)=
  \frac{\mu_0}{4\pi}\,
  \frac{I}{|\bm\ell\times\bm d_A|^2}
  \left(
    \frac{\bm\ell\cdot\bm d_A}{|\bm d_A|}
    -\frac{\bm\ell\cdot\bm d_B}{|\bm d_B|}
  \right)(\bm\ell\times\bm d_A).
  \label{eq:wire}
\end{equation}
Here \(\bm a\) and \(\bm b\) are the segment endpoints in metres,
\(\bm r\) is the field point in metres, \(\bm\ell=\bm b-\bm a\),
\(\bm d_A=\bm r-\bm a\), and \(\bm d_B=\bm r-\bm b\). The current \(I\) is in
amperes and \(\bm F\) is in tesla. The prefactor used by the code is
\(\mu_0/(4\pi)=\SI{1e-7}{\tesla\metre\per\ampere}\), the pre-2019 SI value.
If the field point is within \(\SI{1e-25}{\metre}\) of an endpoint, or if
\(|\bm\ell\times\bm d_A|^2<\SI{1e-70}{\metre\squared}\), the segment returns
zero rather than evaluating the singular expression.

The Johnson--Bovey model replaces the \(\pi\) density by two identical polygonal
loops, one at \(+d_t\hat{\bm n}\) and one at \(-d_t\hat{\bm n}\) from the ring
plane. Writing the \(\si{\angstrom}\)-to-metre conversion implicitly, the
unit-current field at atom position \(\bm r\) for ring type \(t\) is
\begin{equation}
  \bm B^{(1)}_t(\bm r)=
  \sum_i \left[
  \bm F(\bm v_i+d_t\hat{\bm n},\bm v_{i+1}+d_t\hat{\bm n},
         \tfrac{1}{2}\,\si{\nano\ampere},\bm r)
  +\bm F(\bm v_i-d_t\hat{\bm n},\bm v_{i+1}-d_t\hat{\bm n},
         \tfrac{1}{2}\,\si{\nano\ampere},\bm r)
  \right].
  \label{eq:double-loop}
\end{equation}
The code converts \(\si{\angstrom}\) to metres and nanoamperes to amperes at
the boundary of this calculation. The polygon is not a numerical quadrature
mesh; its edges are the physical ring-wire approximation, and each edge is
integrated by the closed form in Eq.~\eqref{eq:wire}.

The tensor kernel stored for that ring and atom is
\begin{equation}
  G^{(t)}_{ab}(\bm r)=-10^6\,\hat n_b\,B^{(1)}_{t,a}(\bm r).
  \label{eq:g-kernel}
\end{equation}
The row index \(a\) is the component of the secondary field, and the column
index \(b\) is the component of the applied field. Equation~\eqref{eq:g-kernel}
is the outer product \(-\bm B^{(1)}_t\hat{\bm n}^{\mathsf T}\), scaled by
\(10^6\). Its trace gives
\begin{equation}
  T_0^{(t)}=-\frac{10^6}{3}\,\hat{\bm n}\cdot\bm B^{(1)}_t .
  \label{eq:bs-t0}
\end{equation}
If a ring-current coefficient \(\kappa_t\) in
\(\si{\nano\ampere\per\tesla}\) is applied later, the usual shielding
contribution would be \(\kappa_t G^{(t)}\) in \(\si{\ppm}\). The present
\texttt{Compute()} path records \(\kappa_t\) for provenance but sums
Eq.~\eqref{eq:g-kernel} with a unit current for every accepted ring.

\begin{center}
\small
\begin{tabular}{lcc}
\hline
Ring type & \(d_t\) (\si{\angstrom}) & \(\kappa_t\) (\si{\nano\ampere\per\tesla})\\
\hline
PHE benzene & 0.64 & -12.00\\
TYR phenol & 0.64 & -11.28\\
TRP benzene, TRP6 & 0.64 & -12.48\\
TRP pyrrole, TRP5 & 0.52 & -6.72\\
TRP indole perimeter, TRP9 & 0.60 & -19.20\\
HIS, HID, HIE imidazole & 0.50 & -5.16\\
\hline
\end{tabular}
\end{center}

For each atom, the calculator queries aromatic rings within
\(\SI{15.0}{\angstrom}\). A candidate is rejected when its center distance is
less than \(\SI{0.1}{\angstrom}\), when that distance is not greater than the
source radius \(0.5(2R)=R\), or when the atom is a ring atom or covalently
bonded to a ring atom. These filters mark where the through-space loop model
is not being used: at singular distances, inside the current distribution, and
in the through-bond region. Accepted ring tensors are summed,
\begin{equation}
  \bm G_i=\sum_{r\in\mathcal A_i}\bm G^{(t_r)}(\bm r_i),
  \label{eq:sum}
\end{equation}
where \(\mathcal A_i\) is the set of accepted nearby aromatic rings for atom
\(i\), and \(t_r\) is ring \(r\)'s type. The sum is decomposed by
Eqs.~\eqref{eq:t0}--\eqref{eq:t2}. The code also accumulates per-ring-type
\(T_0\) and \(T_2\), stores the total unit-current \(\bm B\) field, and records
accepted-ring counts within
\(\SI{3}{\angstrom}\), \(\SI{5}{\angstrom}\), \(\SI{8}{\angstrom}\), and
\(\SI{12}{\angstrom}\) shells.

There is also a free-point sampling path. It evaluates the same unit-current
field and kernel on a grid point, but it does not use atom-topology filters
because a grid point is not a protein atom. It keeps the
\(\SI{0.1}{\angstrom}\) singularity guard, skips points whose center distance
is below the average ring radius, and applies the
\(\SI{15.0}{\angstrom}\) cutoff.

\section{Inputs and output}

The calculator consumes atom coordinates in \(\si{\angstrom}\), typed aromatic
ring topology, ring geometry, and the spatial index used to find nearby rings.
It does not consume MOPAC charges, AIMNet2 charges or embeddings, ff14SB
charges, APBS electrostatic fields, or DSSP secondary-structure labels; those
data feed other calculators. Their internal methods are outside the scope of
this document.

The written arrays are \texttt{bs\_shielding.npy}, an \(N\times9\) packed
\(T_0+T_1+T_2\) kernel; \texttt{bs\_per\_type\_T0.npy}, an \(N\times8\)
per-aromatic-ring-type scalar table; \texttt{bs\_per\_type\_T2.npy}, an
\(N\times40\) table containing eight five-component \(T_2\) blocks;
\texttt{bs\_total\_B.npy}, an \(N\times3\) unit-current magnetic field; and
\texttt{bs\_ring\_counts.npy}, an \(N\times4\) accepted-ring shell-count table.
These arrays are geometric descriptors. Agreement with DFT is assessed by
calibration and correlation, not by expecting the raw kernel to reproduce DFT
shielding pointwise.

\end{document}
